\documentclass[10pt,a4paper]{amsart}
\usepackage[margin=19mm]{geometry}
\usepackage{amsmath,amssymb,mathtools}
\usepackage[scr=rsfs,cal=euler]{mathalfa}
\usepackage{xfrac}
\usepackage[unicode,hidelinks,bookmarks=false]{hyperref}
\newcommand{\ie}{i.e.\ }
\newcommand{\cf}{cf.\ }
\newcommand{\resp}{respectively\ }
\newcommand{\Subsection}{\subsection}
\providecommand{\nobreakdash}{\nobreak\hbox{-}}
\setlength{\emergencystretch}{2em}
\multlinegap0pt
\usepackage{amsthm}
\usepackage{placeins}
\newtheorem{thm}{Theorem}[section]
\newtheorem{exam}{Example}

\newtheorem{lem}[thm]{Lemma}
\newtheorem{prop}[thm]{Proposition}
\newtheorem{cor}[thm]{Corollary}
\makeatletter
\renewcommand{\@seccntformat}[1]{\S{\csname
the#1\endcsname}\hspace{0.5em}}
\makeatother
\theoremstyle{plain}

 
\theoremstyle{definition}
\newtheorem*{dfn}{Definition}
\newtheorem{exa}[thm]{Example}
\newtheorem{rem}[thm]{Remark}
\newcommand{\ds}{\displaystyle}
\newcommand{\bs}{\boldsymbol}
\newcommand{\mb}{\mathbb}
\newcommand{\mc}{\mathcal}
\newcommand{\mf}{\mathfrak}
\renewcommand{\mod}{\operatorname{mod}}
\newcommand{\mult}{\operatorname{Mult}}

\def \a{\alpha} \def \b{\beta} \def \d{\delta} \def \e{\varepsilon} \def \g{\gamma} \def \k{\kappa} \def \l{\lambda} \def \s{\sigma} \def \t{\theta} \def \z{\zeta}

\numberwithin{equation}{section}

\setlength{\leftmargini}{1.em} \setlength{\leftmarginii}{1.em}
\renewcommand{\labelenumi}{\setlength{\labelwidth}{\leftmargin}
   \addtolength{\labelwidth}{-\labelsep}
   \hbox to \labelwidth{\theenumi.\hfill}}
\iffalse
\cite{A}\cite{AA}\cite{BC}\cite{Mag}\cite{C}\cite{Atl}\cite{Da}\cite{GH}\cite{Tot2}\cite{Tot1}\cite{Tot3}
\fi
\title[Strong Gelfand pairs of the symplectic group ${\rm Sp}_4(q)$, $q$ even]{Strong Gelfand pairs of the symplectic group ${\rm Sp}_4(q)$ where $q$ is even: the maximal subgroup $\mathrm{Sp}_2(q^2)\colon 2$ is not a strong Gelfand subgroup (Theorem 5.1)}
\author{Stephen P. Humphries, Joseph E. Marrow}
  \address{Department of Mathematics,  Brigham Young University, Provo, 
UT 84602, U.S.A.
E-mail:  steve@mathematics.byu.edu, jemarrow@math.byu.edu}
\date{Source: arXiv:2504.20279v2 (CC BY 4.0)}
\begin{document}
\begin{abstract}
A strong Gelfand pair $(G,H)$ is a finite group $G$ together with a subgroup $H$ such that every irreducible character of $H$ induces to  a multiplicity-free character of $G$.  We classify the strong Gelfand pairs of the symplectic  groups
${\rm Sp}_4(q)$ for even $q$.


\medskip

\noindent {\bf Keywords}: Strong Gelfand pair, symplectic group, irreducible character, multiplicity one subgroup. \newline 
\medskip
\subjclass[2020]{Primary: 20G40 Secondary: 20C15 20G05}
\end{abstract}
\maketitle
\noindent\emph{Source and licence.} Strong Gelfand pairs of the symplectic group, q even. Version arXiv:2504.20279v2 (CC BY 4.0), distributed under the Creative Commons Attribution 4.0 International licence, http://creativecommons.org/licenses/by/4.0/, as recorded in the arXiv OAI licence metadata for this exact version and shown on the abstract page https://arxiv.org/abs/2504.20279v2. Full rights notice: licenses/arxiv-2504.20279-CC-BY-4.0.txt.\par\smallskip
\noindent\textit{Editorial scope.} Counted unit: the theorem of Section 5 that the pair (Sp\_4(q), Sp\_2(q\^{}2):2) is not a strong Gelfand pair, source0.tex lines 419--421 (an unlabelled thm environment, printed as Theorem 5.1 in the published version and as Theorem 5.1 here), delivered with its complete original author proof at source lines 422--554 as the single primary proof body. No proof marker is added and no mathematics is written, repaired or re-derived. The article is itself an amsart article, so no publisher class is replaced or shipped: the wrapper supplies the pinned runtime preamble, and the article's own declarations are copied verbatim from the source --- the \textbackslash{}newtheorem block at source lines 20--25, the section-number format at source lines 28--31, the theorem styles and the remaining \textbackslash{}newtheorem declarations at source lines 59--65, and the macro block, \textbackslash{}numberwithin and list settings at source lines 67--82. Consequently every printed number of the delivered unit is produced by the article's own counters and coincides with the published version: Sections 1--5 with their unnumbered headings (printed as the section sign followed by the number, as in the published version), Theorem 1.1, Lemmas 2.1--2.2, Propositions 2.3--2.4, Lemma 2.5, Theorem 2.6, Lemma 2.7, Theorem 3.1, Corollary 3.2, Theorem 4.1, Theorem 5.1, the tables and their numbers, and every numbered equation of the delivered sections. Required context retained uncounted: the whole body of the article from the Introduction (source line 86) to the line immediately preceding the counted statement (source line 418), that is the Introduction with the main theorem of the article, the whole of the preliminary section with the total-character lemmas and the maximal-subgroup tables, the whole of Sections 3 and 4 with their theorems and complete proofs, and the opening paragraph of Section 5, together with every definition, lemma, proposition, corollary, table, numbered equation and citation of those source lines. The article's abstract, keywords and subject classification (source lines 47--55) are reproduced verbatim in the wrapper front matter, and the article's own title, author and address lines are reproduced in the wrapper title block. Not part of this unit and not redistributed: the remainder of Section 5 after the counted proof (Corollary 5.3, the case of the maximal subgroup Sp\_4(q\_0), the case of the Suzuki groups and the closing paragraph that concludes the proof of the article's main theorem), the appendix with the Magma listing, the acknowledgement and the article's own bibliography listing. The delivered text cites 21 of the article's bibliography keys; those entries are transcribed verbatim from the article's own inline thebibliography (source lines 650--727) in the article's own order with the article's printed bibliography numbers, so every printed citation number of the delivered unit is the published one, and the one bibliography entry the delivered text does not cite (the Suzuki paper, cited only in the part of Section 5 that lies outside this unit) is not listed. The article writes thirteen of its citations as \textbackslash{}cite followed by a space and then the key list; TeX ignores that space, but the assembler of this corpus scans citation keys with a pattern that does not, so the eleven keys that occur only in that spaced form are declared once more in the wrapper preamble inside an \textbackslash{}iffalse block that TeX never executes and that therefore typesets nothing. That declaration is the only editorial adaptation of this item: it changes no source span, deletes no character, drops no citation, adds no printed character and leaves every printed citation number to be produced by the article's own citation commands and the transcribed bibliography entries, with method exact\_tex for all three delivered bodies and no replacements and no omissions in delivery-evidence.json. No cross-reference of the delivered text leaves the unit: every \textbackslash{}ref and \textbackslash{}eqref of the delivered source lines points at a label defined inside them, so no reference-only adaptation was needed in this item. Printed slips kept literal: the article's preamble loads the youngtab package with the literal option \{...\} (source line 13) and defines an Example environment twice, once unshared at source line 21 and once sharing the theorem counter at source line 64; the delivered unit reproduces the second definition as printed and does not reproduce the unused youngtab load, whose option list is not valid for that package in this runtime; neither is mathematics and neither is visible in the printed unit. The article's own title line contains a misplaced brace in its math group (source line 36); the wrapper title block is editorial and prints the corrected reading, while no source span inside the delivered unit was altered. The counted proof keeps its own typographic forms, including its inline \textbackslash{}qed markers and its running heads. The article contains no figure, no artwork and no \textbackslash{}includegraphics anywhere, so nothing had to be omitted for asset reasons and no artwork is redistributed; the two packages this unit needs beyond the wrapper preamble are amsthm and placeins, the latter for the article's own \textbackslash{}FloatBarrier, and the article's amscd, mathrsfs, graphics, graphicx, caption, epic, enumerate, youngtab, tikz and listings loads are not reproduced because no construct of the delivered source lines uses them. No mathematics is written, repaired or re-derived; all mathematics is native text with no image dependency.
\section{Introduction}



For a finite group $G$ we let $\hat G$ denote the set of irreducible characters of $G$. Then  a {\it multiplicity-free character of $G$} is a character $\chi$ of $G$ such that for   $\psi\in \hat G$, we have $\langle \chi, \psi \rangle \leq 1$. Here only complex characters are considered.
 
A {\it Gelfand pair} $(G,H)$ is a finite group $G$ together with a subgroup $H$ such that the trivial character of $H$ induces a multiplicity-free character of
$G$. The importance of Gelfand pairs is indicated by six equivalent conditions; see     \cite{AHN, BH, Har, Wik}.


A {\it strong Gelfand pair} $(G,H)$ is a finite group $G$ and   $H\le G$ such that for every   $\psi \in \hat H$  the induced character $\psi\uparrow G$ is multiplicity free. 
We will call $H$ a {\it strong Gelfand subgroup of $G$} in this situation. 
Equivalently, $(G,H)$ is a strong Gelfand pair if and only if the Schur ring determined by the {\it $H$-classes} $g^H = \{ g^h : h \in H\}, g\in G$, is commutative \cite{Har, JKar, Den}. Here  our convention is: $g^h=h^{-1}gh$. Note that $(G,G)$ is always a strong Gelfand pair.

In this paper we continue our investigation of strong Gelfand pairs of  groups that are close to being simple; in \cite {AHN, BH, GH}  we found all such pairs for 
 $G={\rm SL}(2,p^n),n\ge 1,$  $p$ a prime, and  the  symmetric groups. We refer to \cite {BH,Har, GH} for necessary background and to \cite {BC} for some of  the latest results on strong Gelfand pairs.
 

 We note that Gelfand pairs and strong Gelfand pairs have applications in representation theory; see  \cite {A,AA,BC,C} among many other references.
As explained above, an equivalent condition for  $(G,H)$ to be  a strong Gelfand pair is that the Schur ring determined by the $H$-classes is commutative.
This shows that a strong Gelfand pair determines a commutative Schur ring and so a commutative  association scheme, which then gives  indicates connections with algebraic combinatorics.
One other application of strong Gelfand pairs is to random walks on finite groups: if $(G,H)$ is  a strong Gelfand pair, then one can define a random walk on $G$ using probabilities that are constant on the above mentioned $H$-classes. The commutativity property of the $H$-classes means that the random walk is `diagonalizable' and so can be very well 
understood.


This paper will consider strong Gelfand pairs for the symplectic groups {\rm Sp}$_4(2^n)$ as their irreducible characters are known \cite{speven}. In contrast, the irreducible characters of ${\rm Sp}_{2k}(q),k>2,$ are not understood, where $q$ is a prime power. The  groups {\rm Sp}$_4(2^n), n>1,$ are simple and  the representation theory of these groups is considered in \cite {Da}.
 The main result of this paper is: 

\begin{thm}\label{Thm1}

The only strong Gelfand pair {\rm (Sp}$_4(2^n),H),n\ge 2,$ is where $H={\rm Sp}_4(2^n)$. 
\end{thm}
 

Throughout we will use the standard  Atlas notation \cite {Atl}.
\medskip 




\section{Preliminary results}
\FloatBarrier


All groups considered in this paper will be assumed finite. 
For a  group $G$, the \textit{total character} of $G$, denoted $\tau_G$, is the sum of all the irreducible characters of $G$; see \cite {Tot1,Tot3,Tot2}. The following gives the `total character argument' for showing that certain subgroups are not strong Gelfand subgroups:

\begin{lem} [Lemma $3.3$ \cite {GH}] \label{lem:totalchar} Let $H\leq G$ be   groups. If there is  $\chi\in \hat G$ with 
$$
\deg(\tau_H) < \deg(\chi),
$$
 then $(G, H)$ is not a strong Gelfand pair.
\end{lem}

The following indicates that it is important to determine which maximal subgroups are strong Gelfand pairs. 


\begin{lem}  [Lemma $3.1$ \cite {BH}]  \label{lem:totalstack} Suppose we have groups $H\leq K\leq G$. If $(G, K)$ is not a strong Gelfand pair, then neither is $(G, H)$.
\end{lem}



For  $q=2^e$,  $e>1$, we find from Table $8.14$ of \cite{max} that   the maximal subgroups of $\mathrm{Sp}_4(q)$  are as listed in Table \ref{maxeven}.

\begin{table}[h!]
\centering
\caption{Maximal subgroups of $\mathrm{Sp}_4(q)$, for $q=2^e, e>1$}
\label{maxeven}
\vspace{5pt}
\begin{tabular}{|ccc|}
\hline
Group & Order& Conditions\\
\hline
$\vphantom{E^{3^2}}E_q^3\colon\mathrm{GL}_2(q)$ & $q^3\cdot (q^2+q)(q-1)^2$ & \\%
$\vphantom{E^{3^2}}E_q^3\colon\mathrm{GL}_2(q)$ & $q^3\cdot (q^2+q)(q-1)^2$ & \\%
$\mathrm{Sp}_2(q)\wr 2$ & $2q^2(q^2-1)^2$ & \\%
$\mathrm{Sp}_2(q^2)\colon 2$ & $2q^2(q^4-1)$ & \\
$\mathrm{Sp}_4(q_0)$ & $q_0^4(q_0^2-1)(q_0^4-1)$ & $q=q_0^r, r$ is prime\\%
$\mathrm{SO}_4^+(q)$ & $2q^2(q^2-1)^2$ & \\
$\mathrm{SO}_4^-(q)$ & $2q^2(q^4-1)$ & \\
$\mathrm{Sz}(q)$ & $q^2(q^2+1)(q-1)$ & $e$ odd\\%
\hline
\end{tabular}
\end{table}

Table \ref{maxeven} has the maximal  subgroup $E_q^3 \colon \mathrm{GL}_2(q)$ listed twice because there are two  non-conjugate maximal subgroups of Sp$_4(q)$ which are isomorphic to $E_q^3 \colon \mathrm{GL}_2(q)$. 


From Lemma \ref {lem:totalstack} Theorem \ref {Thm1} will follow if we can show that none of these maximal subgroups is a strong Gelfand subgroup. We consider each case separately.

The next two results  will allow us to assume   $e\ge 3$.   

We first  consider  the symplectic group ${\rm Sp}_4(2)$; since ${\rm Sp}_4(2)\cong S_6$ the result here follows from our consideration of the symmetric groups in 
\cite {AHN}:

\begin{prop} \label{thm:sp2}
The only proper subgroups of $\mathrm{Sp}_4(2)$ which are strong Gelfand subgroups are  the maximal subgroups.\qed
\end{prop}





\begin{prop}\label{cor:sp4}
No proper subgroup of $\mathrm{Sp}_4(4)$ is a strong Gelfand subgroup.
\end{prop}
\noindent {\it Proof} 
 We use the MAGMA \cite {Mag} code given in the Appendix  to obtain  this result.\qed\medskip
 
 In what follows we will often have the situation where $H \le G, |G:H|=2$. We introduce the following conventions.
 For $\psi \in \hat H$ it  is well-known \cite {JaL} that either 
 
 \noindent (i) $\psi \uparrow G$ is a sum of two distinct characters  in $\hat G$  (call this the {\it splitting case}); or
 
 \noindent (ii)  $\psi \uparrow G$ is irreducible (call this the {\it fusion} case). 
 
 In the splitting case, if 
  $\psi \uparrow G=\chi_1+\chi_2, \chi_1,\chi_2 \in \hat G$, then $\chi_i\downarrow H=\psi, i=1,2.$
 
\medskip

In the situation $|G:H|=2$ the relationship between $\tau_G$ and $\tau_H$ is given in: 
\begin{lem}\label{lemsplitf}
Let $H\le G, |G:H|=2$. Let $\mathcal S$ be the set of $\psi \in\hat H$ that split and 
let  $\mathcal F$ be the set of $\psi \in\hat H$ that fuse. Then 
$$
\tau_G(1)=2\sum_{\psi \in \mathcal S}\psi(1)+  \sum_{\psi \in \mathcal F}\psi(1).\qed 
$$
\end{lem}



The character table for $\mathrm{Sp}_4(q)$ is given in \cite{speven} and  we will use notation from \cite{speven}. 

\begin{thm}  \cite{speven}.\label{lem:totalsymplectic}  (i) The degree of the total character of $\mathrm{Sp}_4(q)$ is $q^6+q^4-q^2$ if $q$ is even.

(ii) The largest degree of an irreducible character of $\mathrm{Sp}_4(q)$ is $q^4+2q^3+2q^2+2q+1$ when $q\geq 4$ is a power of $2$.
\end{thm}
\noindent\textit{Proof.} (i) We just  sum  the degrees of characters of $\mathrm{Sp}_4(q)$ as listed in \cite{speven}. (ii) follows directly from \cite {speven}.\qed\medskip


\begin{lem}\label{isomorpicmaximal}
If $q=2^e, e>1$, then $\mathrm{Sp}_2(q)\wr 2 \cong \mathrm{SO}_4^+(q)$  and $\mathrm{Sp}_2(q^2)\colon 2 \cong \mathrm{SO}_4^-(q)$.
\end{lem}
\noindent\textit{Proof.} 
See Proposition $7.2.1$  and Table $8.14$ of \cite{max}.  \qed\medskip 


We now consider the maximal subgroups separately in the following sections.

\section { The maximal subgroup $\mathrm{Sp}_2(q)\wr 2 $}

%xxxx

\begin{thm}\label{wreath}
For $q=2^e, e>1,$ the maximal subgroup $\mathrm{Sp}_2(q)\wr 2 \leq \mathrm{Sp}_4(q)$ is not a strong Gelfand subgroup.
\end{thm}
\noindent\textit{Proof.} This proof will be a `total character argument' and so we will need to find the total character of
$\mathrm{Sp}_2(q)\wr 2$. 
We have $\mathrm{Sp}_2(q)\wr 2 = \mathrm{Sp}_2(q)^2: 2$ and one way to  represent the elements of $\mathrm{Sp}_2(q)^2: 2$ is by $2 \times 2$ blocks of $2 \times 2$ matrices, where the cyclic  subgroup $2$ is generated  by  $\begin{bmatrix} 0&I_2\\I_2&0\end{bmatrix}$ and
$(a,b) \in \mathrm{Sp}_2(q)^2$ is represented as the block matrix 
$\begin{bmatrix} a&0\\0&b\end{bmatrix}.$ 
 

Now $\mathrm{Sp}_{2}(q) \cong \mathrm{SL}_2(q)$ has  character table  given in \cite{Dor}  (see also  \cite {GH}); we reproduce it here in Table \ref{sleven}.
Here the parameters $s,t,j,m$ satisfy $1\leq s, t \leq (q-2)/2$, $1\leq j, m \leq q/2$,  $\rho$ is a primitive $(q-1)$-th root of unity and $\sigma$ a primitive $(q+1)$-th root of unity. 

\begin{table}[h!]
\centering
\caption{Character Table for $\mathrm{SL}_{2}(q)$ with $q$ even}
\label{sleven}
\vspace{5pt}
\begin{tabular}{c|cccc}
Class & $1$ & $c$ & $a^t$ & $b^m$\\
Size & $1$ & $q^2-1$ & $q(q+1)$ & $q(q-1)$\\
\hline
 Tr & $1$ & $1$ & $1$ & $1$\\
$\psi$ & $q$ & $0$ & $1$ & $-1$\\
$\chi_s$ & $q+1$ & $1$ & $\rho^{st}+\rho^{-st}$ & $0$\\
$\theta_j$ & $q-1$ & $-1$ & $0$ & $-(\sigma^{jm}+\sigma^{-jm})$\\
\end{tabular}
\end{table}

Here the conjugacy classes of SL$_2(q)$ are represented by powers of the following elements:
$$
1 = \begin{bmatrix}
1 & 0\\
0 & 1
\end{bmatrix},\ 
c=\begin{bmatrix}
1 & 0\\
1 & 1
\end{bmatrix},\ 
a=\begin{bmatrix}
\rho & 0\\
0 & \rho^{-1}
\end{bmatrix},
$$
and an element $b$ of order $q+1$. We also give the sizes of the classes in Table \ref{sleven}.

Since $\mathrm{Sp}_2(q)\wr 2 \cong \mathrm{Sp}_2(q)^2\colon 2$
  the irreducible characters of $\mathrm{Sp}_2(q)\wr 2$ are easily found using Table 2.
In Table 
   3 we give the degrees of the irreducible characters of $\mathrm{Sp}_2(q)\wr 2$.
These character degrees are obtained using \cite[Proposition $20.9$, Theorem $19.18$]{JaL}. Further, in Table 
 3 we are assuming that  
 $$1 \leq s, s^\prime \leq (q-2)/2, \,\,1 \leq j, j^\prime \leq q/2 \text { and } s\neq s^\prime,j\neq j^\prime.
 $$ 



\begin{table}[h!]
\label{chardegwreath}
\centering
\caption{Character degrees for $\mathrm{Sp}_2(q)\wr 2$ with $q$ even}
\vspace{5pt}
\begin{tabular}{c|cc}
Character & Degree & Multiplicity\\
\hline
$\vphantom{()^{2^2}}(\mathrm{ Tr}\times\mathrm{ Tr})_1$ & $1$ & $1$\\
$(\mathrm{ Tr}\times\mathrm{ Tr})_2$ & $1$ & $1$\\
$\mathrm{ Tr}\times\psi$ & $2\cdot q$ & $1$\\
$\mathrm{ Tr}\times\chi_i$ & $2\cdot(q+1)$ & $(q-2)/2$\\
$\mathrm{ Tr}\times\theta_j$ & $2\cdot (q-1)$ & $q/2$\\
$(\psi\times\psi)_1$ & $q^2$ & $1$\\
$(\psi\times\psi)_2$ & $q^2$ & $1$\\
$\psi\times\chi_s$ & $2\cdot q(q+1)$ & $(q-2)/2$\\
$\psi\times\theta_j$ & $2\cdot q(q-1)$ & $q/2$\\
$(\chi_s\times\chi_s)_1$ & $(q+1)^2$ & $(q-2)/2$\\
$(\chi_s\times\chi_s)_2$ & $(q+1)^2$ & $(q-2)/2$\\
$\chi_s\times\chi_{s^\prime}$ & $2\cdot (q+1)^2$ & $(q-2)(q-4)/8$\\
$\chi_s\times\theta_j$ & $2\cdot(q^2-1)$ & $q(q-2)/4$\\
$(\theta_j\times\theta_j)_1$ & $(q-1)^2$ & $q/2$\\
$(\theta_j\times\theta_j)_2$ & $(q-1)^2$ & $q/2$\\
$\theta_j\times\theta_{j^\prime}$ & $2\cdot (q-1)^2$ & $q(q-2)/8$
\end{tabular}
\end{table}

In Table 3 the suffices $1,2$ are written to indicate that these  are the split cases. The lack of such a suffix indicates the fusion cases. 
In Table 3 each case has a certain `Multiplicity' that is also indicated; this  depends on the parameters involved. 
Then from 
 Table 3 we obtain the degree of the total character of  Sp$_2(q)\wr 2$:
\begin{align*}
(\tau_{\mathrm{Sp}_2(q)\wr 2})(1)
&=1 + 1 + 2 q + (q + 1) (q - 2) + (q - 1) \cdot q + 2 q^2 + q\cdot(q + 1)\cdot (q - 2)\\ & \qquad + q^2 (q - 1)
 + (q + 1)^2 \cdot(q - 2) + (q + 1)^2\cdot (q - 2)\cdot (q - 4)/4\\& \qquad + (q^2 - 1)\cdot q\cdot (q - 2)/2 + (q - 1)^2\cdot q
 + (q - 1)^2\cdot q\cdot (q - 2)/4 \\&
= q^4+q^3-q.
\end{align*}



Now $q^4+q^3-q < q^4+2q^3+2q^2+2q+1$, and  by Theorem  \ref{lem:totalsymplectic}   $q^4+2q^3+2q^2+2q+1$ is the degree of an irreducible character of $\mathrm{Sp}_4(q)$. Then by Lemma \ref{lem:totalchar}  $(\mathrm{Sp}_4(q), \mathrm{Sp}_2(q)\wr 2)$ is not a strong Gelfand pair.\qed\medskip 

By Lemma \ref {isomorpicmaximal} and the fact that the above argument is a `total character argument' (not dependent on the particular embedding of $ \mathrm{Sp}_2(q)\wr 2$ in $\mathrm{Sp}_4(q)$)        we see that we have now also dealt with  the maximal subgroup ${\rm SO}^+_4(q)$ case from Table 1:

\begin{cor} \label{cor1} The maximal subgroup {\rm SO}$^+_4(q)<{\rm Sp}_4(q)$ is not a strong Gelfand subgroup.\qed
\end{cor}\medskip
 


\section {The maximal subgroups $E_q^3:{\rm GL}_2(q)$}

By Theorems \ref {thm:sp2} and \ref {cor:sp4}  we may assume that $q>4$.
\begin{thm}\label{thm:elementaryabelian}
For $q=2^e, e>2$, the maximal subgroup $E_q^3\colon\mathrm{GL}_2(q) \leq \mathrm{Sp}_4(q)$ is not a strong Gelfand subgroup.
\end{thm}
\noindent\textit{Proof.} In \cite{speven} two isomorphic maximal subgroups are considered; they are denoted $P$ and $Q$. 
The orders of $P$ and $Q$ are $q^3(q^2+q)(q-1)^2$ and they are isomorphic to  $E_q^3\colon\mathrm{GL}_2(q)$. The character tables for these subgroups are given in \cite{speven}. 


We  take the inner product of the character of $P$ denoted by  $\chi_5(k)$ in  \cite{speven} with a character of $\mathrm{Sp}_4(q)$ restricted to $P$, namely $\chi_{1}(m, n)\downarrow P$. In what follows  $A_i,A_{ij}, C_j, D_k$ is the notation used  in \cite{speven} for the classes of $P$; further, the sizes of these classes are also given in \cite{speven}. Using all of this information we obtain:
\begin{align*}
&\langle \chi_5(k), \chi_1(m, n)\downarrow P\rangle\\&
=\frac{1}{|P|}\Big(
|A_1| q(q^2-1) (q+1)^2(q^2+1)
+|A_2| q(q-1) (q+1)^2 
+|A_{31}| (-q)(q+1) (q+1)^2\\
&\qquad + |A_{32}| (-q) (2q+1)
+|C_2(i)| (q-1)\alpha_{ik} (q+1)\alpha_{im}\alpha_{in} +|D_2(j)| (-\alpha_{jk}) \alpha_{jm}\alpha_{jn}
\Big)\\&
=\frac{1}{q^4(q-1)(q^2-1)}\Big(q(q^2-1) (q+1)^2(q^2+1)\\
&\qquad +(q^2-1) q(q-1) (q+1)^2
+(q-1) (-q)(q+1) (q+1)^2\\ 
&\qquad +(q-1)(q^2-1) (-q) (2q+1)\\ 
&\qquad +\sum_{i=1}^{(q-2)/2} q^3(q+1) (q-1)\alpha_{ik} (q+1)\alpha_{im}\alpha_{in}
+\sum_{j=1}^{(q-2)/2} q^3(q^2-1) (-\alpha_{jk}) \alpha_{jm}\alpha_{jn}\Big)\\
&=\frac{1}{q^7-q^6-q^5+q^4}\Big(q^7+2q^6+q^5-q^3-2q^2-q
 +q^6+q^5-2q^4-2q^3+q^2\\
 &\qquad \qquad \qquad +q
-q^5-2q^4+2q^2+q-2q^5+q^4+3q^3-q^2-q\\
&\qquad +\sum_{i=1}^{(q-2)/2} (q^6+q^5-q^4-q^3)\alpha_{ik}\alpha_{im}\alpha_{in}
+\sum_{j=1}^{(q-2)/2}(-q^5+q^3)\alpha_{jk}\alpha_{jm}\alpha_{jn}\Big)\\
\end{align*}
\begin{align*}&=
\frac{1}{q^7-q^6-q^5+q^4} \left(q^7+3q^6-q^5-3q^4+\left(q^6-q^4\right)\sum_{j=1}^{(q-2)/2} \alpha_{jk}\alpha_{jm}\alpha_{jn}\right)\\&
=\frac{3+q}{q-1}+\frac{1}{q-1}\left(\sum_{j=1}^{(q-2)/2} \alpha_{jk}\alpha_{jm}\alpha_{jn}\right).
\end{align*}

Here $\alpha_{ij} = \overline{\gamma}^{ij}+\overline{\gamma}^{-ij}$ where $\langle\gamma\rangle = \mathbb{F}_q^\times$, and $\overline{\gamma}$ is the image of $\gamma$ under a fixed monomorphism from $\mathbb{F}_q^\times$ into $\mathbb{C}^\times$, making $\overline{\gamma}$ a $(q-1)$-th root of unity. For clarity of notation, in our calculations we omit the overline.

Now supposing that $q>5$, if we have $\sum_{i=1}^{(q-2)/2}\alpha_{ik}\alpha_{im}\alpha_{in}=q-5$, then the above gives $\langle \chi_5(k), \chi_1(m, n)\downarrow P\rangle=2$. We will now show that there is a choice of $k, m, n$ so that $\sum_{i=1}^{(q-2)/2}\alpha_{ik}\alpha_{im}\alpha_{in}$ is equal to  $q-5$.
We calculate:
\begin{align*}
&\sum_{j=1}^{(q-2)/2}\alpha_{jk}\alpha_{jm}\alpha_{jn}\\
&= \left(\sum_{j=1}^{(q-2)/2} \gamma^{jk}+\gamma^{-jk}\right)\left(\sum_{j=1}^{(q-2)/2} \gamma^{jm}+\gamma^{-jm}\right)\left(\sum_{j=1}^{(q-2)/2} \gamma^{jn}+\gamma^{-jn}\right)\\
&=\sum_{j=1}^{q-2} \gamma^{j(k+m+n)} + \sum_{j=1}^{q-2} \gamma^{j(k+m-n)} + \sum_{j=1}^{q-2} \gamma^{j(k-m+n)} + \sum_{j=1}^{q-2} \gamma^{j(k-m-n)}
\end{align*}
and notice that each of these four sums will be $q-2$ if $q-1$ divides  $j$, and $-1$ otherwise. Suppose that $q>5$ and choose $k=q-4, m=1, n=2$. Then  $m\neq n$ and $m+n\neq q-1$, as required. We also have that only one of $k\pm m \pm n $ is congruent to zero mod
  ${q-1}$. This gives
\begin{align*}
&\sum_{j=1}^{q-2} \gamma^{j(q-1)} + \sum_{j=1}^{q-2} \gamma^{j(q-3)} + \sum_{j=1}^{q-2} \gamma^{j(q-5)} + \sum_{j=1}^{q-2} \gamma^{j(q-7)} = q-5.
\end{align*}

Then for $k=q-4, m=1, n=2$ we have:
\begin{align*}
&\langle \chi_5(q-4), \chi_1(1, 2)\downarrow P\rangle = \frac{3+q+\left(\sum_{j=1}^{(q-2)/2} \alpha_{j(q-4)}\alpha_{j1}\alpha_{j2}\right)}{q-1}
 = \frac{3+q+q-5}{q-1}
= 2,
\end{align*}
showing that $(\mathrm{Sp}_4(q), P)$ is not a strong Gelfand pair if $q>5$.

A similar argument shows that $(\mathrm{Sp}_4(q), Q), q>5,$ is also not a strong Gelfand pair.
\qed
\medskip



\section {The maximal subgroups Sp$_2(q^2):2$ and {\rm Sp}$_4(q_0)$}

The elements of  the field $\mathbb F_{q^2}$ can be represented as $2 \times 2$ matrices over $\mathbb F_q$. This shows how Sp$_2(q^2)\le
{\rm Sp}_4(q)$. The action of the $2$ in Sp$_2(q^2):2$ is the Galois action. 


\begin{thm}
For $q=2^e, e>1,$ the pair $\left(\mathrm{Sp}_4(q), \mathrm{Sp}_2\left(q^2\right)\colon 2\right)$ is not a strong Gelfand pair.
\end{thm}

\noindent\textit{Proof.} Let $G=\mathrm{Sp}_2\left(q^2\right)\colon 2$ and $H = \mathrm{Sp}_2\left(q^2\right)\leq G$. Using  Table \ref{sleven} we get the character table for $H$; see  Table \ref{lineardegrees} where  $1 \leq s \leq (q^2-2)/2$ and $1 \leq j \leq q^2/2$.

\begin{table}[h!]
\centering
\caption{Character degrees for $\mathrm{Sp}_{2}\left(q^2\right)$ with $q$ even}
\vspace{5pt}
\label{lineardegrees}
\begin{tabular}{c|c c}
Character & Degree&Multiplicity\\
\hline
 Tr & $1$& $1$\\
$\psi$ & $q^2$& $1$\\
$\chi_s$ & $q^2+1$ & $(q^2-2)/2$\\
$\theta_j$ & $q^2-1$&$q^2/2$
\end{tabular}
\end{table}

In order to find the degree of $\tau_{G}$, we will need to determine which characters of $H$ split and which  fuse; 
it will suffice to determine which characters of $H$ induce to  irreducible characters of $G$. Again from \cite{JaL}, since $|G \colon H|=2,$ we know that, by  inducing, every character in $\hat{H}$ either splits into a sum of two irreducible characters or fuses pairwise into irreducible characters in $\hat{G}$.
We  use Lemma \ref {lemsplitf} and Table 4 to give:

 
\begin{prop}\label{propsplit} Let $G={\rm Sp}_2(q^2):2\ge H={\rm Sp}_2(q^2)$. Then 

\noindent (i) $\text{Tr}_H$ splits;

\noindent (ii)  $\psi$ splits;

\noindent (iii) all $\theta_j$ fuse;

 The characters $\chi_s$ sometimes split, but not always: 
 
\noindent (iv)  $\chi_s\uparrow G$ is  irreducible  if $(q^2-1)\nmid s(q\pm1)$; and 
 
 
 \noindent (v) $\chi_s\uparrow G$ is the sum of two irreducible characters if $(q^2-1) \mid s(q\pm1)$.
\end{prop}
\noindent{\it Proof} (i) It is clear that Tr$_H$ splits.

\noindent (ii) Since $\psi(1)=q^2$ and there is no other character of degree $q^2$ we see that $\psi$ cannot fuse.

\noindent (iii)  It will suffice to show that $\langle \theta_j \uparrow G,\theta_j \uparrow G\rangle=1$. 
Now a calculation shows that $\theta_j \uparrow \mathrm{Sp}_2(q^2)\colon 2$ 
is as  described in the following  table, where $\sigma$ is a primitive $(q^2+1)$-th root of unity. 
\\[5mm]
\begin{tabular}{c|ccccc}
 & $\text{Tr}_H$ & $c$ & $a^t$ & $b^m$ & $G \smallsetminus H$\\
\hline\\
$\theta_j \uparrow \mathrm{Sp}_2(q^2)\colon 2$ & $2q^2-2$ & $-2$ & $0$ & $-(\sigma^{jm} + \sigma^{-jm}+\sigma^{jmq} + \sigma^{-jmq})$ & $0$
\end{tabular}
\\[5mm]
Now 
$(\theta_j\uparrow G)(G \setminus H)=\{0\}$ and 
for $g \in H$ we have $g$ and $g^{-1}$ are conjugate. Thus 
\begin{align*}
\langle \theta_j\uparrow G, \theta_j\uparrow G\rangle &= \frac{1}{|G|} \sum_{g\in G}(\theta_j\uparrow G)(g)\cdot (\theta_j\uparrow G)\left(g^{-1}\right)\\&=
\frac{1}{|G|} \sum_{g\in H}(\theta_j\uparrow G)(g)\cdot (\theta_j\uparrow G)(g^{-1})
=\frac{1}{|G|} \sum_{g\in H}(\theta_j\uparrow G)^2(g).
\end{align*}

Using Table 2 again and taking $g_m \in (b^m)^G$  the above is equal to
\begin{align*}&
\frac{1}{2(q^6-q^2)}\left(\underbrace{(2q^2-2)^2}_{\text{ Tr}} + \underbrace{(-2)^2(q^4-1)}_{c} + \underbrace{0}_{\text{$a^t$}}+\underbrace{(q^4-q^2)}_{\text{size of $(b^m)^H$}} {\sum_{m=1}^{q^2/2}}\left(\theta_j\uparrow G\right)^2(g_m)\right)
\\&=
\frac{1}{2(q^6-q^2)}\left({8q^4-8q^2} + {(q^4-q^2)} {\sum_{m=1}^{q^2/2}}  {\left(\theta_j\uparrow G\right)^2(g_m)}\right)
\\&=
\frac{1}{2(q^6-q^2)}\left(8q^4-8q^2+(q^4-q^2)\sum_{m=1}^{q^2/2}\left(-\sigma^{jm}-\sigma^{-jm}-\sigma^{jmq}-\sigma^{-jmq}\right)^2\right)
\\&=
\frac{1}{2(q^6-q^2)}
\bigg (
8q^4-8q^2+(q^4-q^2)\sum_{m=1}^{q^2/2}
\big (
4 +  {(\sigma^{2jm}+\sigma^{-2jm})}+ (\sigma^{2jmq} +\sigma^{-2jmq})\\&\qquad  + (2\sigma^{jm(q+1)} + 2\sigma^{-jm(q+1)}) + (2\sigma^{jm(q-1)} + 2\sigma^{-jm(q-1)})
\big )
\bigg )
\end{align*}
 

Now, since $1+\sum_{i=1}^{q^2/2} \sigma^i+\sigma^{-i}=\sum_{i=0}^{q^2} \sigma^i$,  the above  is
\begin{align*}&
\frac{1}{2(q^6-q^2)}\left(8q^4-8q^2+(q^4-q^2)\sum_{m=1}^{q^2}\left(
2+\sigma^{2jm}+\sigma^{2jmq}+2\sigma^{jm(q+1)}+2\sigma^{jm(q-1)}
\right)\right)\\&\qquad = 
\frac{(8q^4-8q^2+2q^2(q^4-q^2)-6(q^4-q^2))}{2q^6-2q^2} = \frac{2q^6-2q^2}{2q^6-2q^2} = 1
\end{align*}
as required for (iii). 
\medskip


\noindent (iv) 
Now a calculation shows that $\chi_s \uparrow \mathrm{Sp}_2(q^2)\colon 2$ 
is as  described in the following  table, where $\rho$ is a primitive $(q^2-1)$-th root of unity. 
\\[5mm]
\begin{tabular}{c|ccccc}
 & $\text{Tr}_H$ & $c$ & $a^t$ & $b^m$ & $G \smallsetminus H$\\
\hline\\
$\theta_j \uparrow \mathrm{Sp}_2(q^2)\colon 2$ & $2q^2+2$ & $2$ &  $-(\rho^{jm} + \rho^{-jm}+\rho^{jmq} + \rho^{-jmq})$ &$0$  & $0$
\end{tabular}
\\[5mm]
We again examine $\langle \chi_s,\chi_s\rangle$ to see when we obtain $1$.
 Taking $g_t\in (a^t)^G$ an argument similar to the $\theta_j$ case gives   
 \begin{align*}
 &\langle \chi_s\uparrow G, \chi_s  \uparrow G\rangle
 =\frac{1}{|G|} \sum_{g\in G} \left(\chi_s\uparrow G\right)(g)\cdot\left(\chi_s\uparrow G\right)\left(g^{-1}\right)\\
 =&\frac{1}{2q^6-2q^2}\left(\sum_{g\in H} \left(\chi_s\uparrow G\right)^2(g)\right)\\
 =&\frac{1}{2q^6-2q^2}\left((2q^2+2)^2 + 4(q^4-1) + (q^4+q^2)\sum_{t=1}^{(q^2-2)/2}\left(\chi_s(g_t)\right)^2\right)\\
 =&\frac{1}{2q^6-2q^2}\left(8q^4+8q^2 + (q^4+q^2)\sum_{t=1}^{(q^2-2)/2}\left(\rho^{st} + \rho^{-st} + \rho^{stq} + \rho^{-stq}\right)^2\right)\\
 =&\frac{1}{2q^6-2q^2}\Bigg(8q^4+8q^2 + (q^4+q^2)\sum_{t=1}^{(q^2-2)/2}\Big(4+\rho^{2st} + \rho^{-2st} + \rho^{2stq} + \rho^{-2stq}\\
  &\qquad \qquad\qquad\qquad\qquad\qquad\qquad\qquad\qquad + 2\rho^{-st(q-1)} + 2\rho^{-st(q+1)}\Big)\Bigg)\\
 =& \frac{1}{2q^6-2q^2}\left(8q^4+8q^2 + \left(2(q^2-2)-2\right)(q^4+q^2) + (q^4+q^2)\sum_{t=1}^{q^2-2}(2\rho^{st(q+1)} + 2\rho^{st(q-1)})\right)\\
 =& \frac{1}{2q^6-2q^2}\left(8q^4+8q^2 + 2q^6-4q^4-6q^2 +  (q^4+q^2)\sum_{t=1}^{q^2-2}(2\rho^{st(q+1)} + 2\rho^{st(q-1)})\right)\\
 =& \frac{1}{2q^6-2q^2}\left(2q^6+4q^4+2q^2 +  (q^4+q^2)\sum_{t=1}^{q^2-2}(2\rho^{st(q+1)} + 2\rho^{st(q-1)})\right).
 \end{align*}

Here we used the facts that $(q^2-1)\nmid 2s$ and $(q^2-1) \nmid 2sq$, since $1\le s\le (q^2-2)/2$. Now, since $(q^2-1) \nmid s$, only one of $(q^2-1) \mid s(q+1)$ or $(q^2-1) \mid s(q-1)$ can be true, this shows that the above is equal to 
$$ \begin{cases}
\frac{1}{2q^6-2q^2}\left(2q^6+4q^4+2q^2 -4(q^4+q^2)\right) = \frac{2q^6-2q^2}{2q^6-2q^2} =1 \text{ if } (q^2-1) \nmid s(q\pm 1)\\
\frac{1}{2q^6-2q^2}\left(2q^6+4q^4+2q^2 +  2(q^4+q^2)(q^2-3)\right) = \frac{4q^6-4q^2}{2q^6-2q^2} = 2 \text{ if } (q^2-1) \mid s(q\pm 1).
\end{cases}
$$

Since there are $\frac{q}{2}$ values of $s$ for which $(q^2-1) \mid s(q+1)$ and $\frac{q-2}{2}$ values where $(q^2-1) \mid s(q-1)$,  we see  that $\frac{2q-2}{2}=q-1$ characters $\chi_s$ of $H$ split in $G$. Then the remaining $\frac{q^2-2q}{2}$ characters fuse in $G$. Recall that $1\leq j \leq q^2/2$ and $1 \leq s \leq (q^2-2)/2$. 
So 
$$
\deg\left(\tau_{G}\right) = 2 + 2q^2 + \left(2(q-1) + \frac{q^2-2q}{2}\right)(q^2+1) + \frac{q^2}{2}\left(q^2-1\right) = q^4+q^3+q.
$$

By Theorem  \ref {lem:totalsymplectic}  $q^4+2q^3+2q^2+2q+1$ is the largest degree of an irreducible character of $\mathrm{Sp}_4(q), q\geq 4$.   Since $G = \mathrm{Sp}_2\left(q^2\right)\colon 2$ and
$$
\deg\left(\tau_{G}\right) =  q^4+q^3+q < q^4+2q^3+2q^2+2q+1
$$
by Lemma \ref {lem:totalchar}  $\left(\mathrm{Sp}_4(q), \mathrm{Sp}_2\left(q^2\right)\colon 2\right)$ is not a strong Gelfand pair.
\qed

\begin{thebibliography}{99}
\bibitem[1]{A} Aizenbud, A.; Gourevitch, D.; Rallis, S.; Schiffmann, Gerard G., \emph{Multiplicity one theorems}, Ann. Math. (2) 172(2) (2010) 1407--1434.
\bibitem[2]{AA} Aizenbud, A.; Gourevitch, D., \emph{Multiplicity one theorem for $({\rm GL}_{n+1} (\mathbb R), {\rm GL}_n (\mathbb R))$}, Selecta Math. (N.S.) 15(2) (2009) 271--294.
\bibitem[3]{AHN} Anderson, Gradin; Humphries, Stephen P.; Nicholson, Nathan \emph{Strong Gelfand pairs of symmetric groups.} J. Algebra Appl. 20 (2021), no. 4, Paper No. 2150054, 22 pp.
\bibitem[4]{BH} Barton, Andrea; Humphries, Stephen \emph{Strong Gelfand Pairs of {\rm SL}($2,p$)}. J. Algebra Appl. 22 (2023), no. 6, Paper No. 2350133, 13 pp. https://doi.org/10.1142/S0219498823501335
\bibitem[5]{max} Bray, John N.; Holt, Derek F.; Roney-Dougal, Colva M. \emph{The maximal subgroups of the low-dimensional finite classical groups. With a foreword by Martin Liebeck.} London Mathematical Society Lecture Note Series, 407. Cambridge University Press, Cambridge, 2013. MR3098485
\bibitem[6]{BC} Brou, Kouakou Germain; Coulibaly, Pie; Kangni, Kinvi, \emph{Generalized Gabor transform for a strong Gelfand pair.} J. Adv. Math. Stud. 18 (2025), no. 1, 109--121.
\bibitem[7]{Mag} Bosma, Wieb; Cannon, John; Playoust, Catherine; \emph{The Magma algebra system. I. The user language}, J. Symbolic Comput., 24 (1997), 235--265.
\bibitem[8]{C} Chan, Kei Yuen, \emph{Ext-multiplicity theorem for standard representations of $({\rm GL}_{n+1},{\rm GL}_{n})$.} Math. Z. 303 (2023), no. 2, Paper No. 45, 25 pp.
\bibitem[9]{Har} Ceccherini-Silberstein, T.; Scarabotti, F.; Tolli, F., \emph{Harmonic analysis on finite groups, in Representation Theory, Gelfand Pairs and Markov Chains}, Vol. 108 Cambridge Studies in Advanced Mathematics, Cambridge University Press, Cambridge, 2008, xiv + 440 pp.
\bibitem[10]{Atl} Conway, J. H.; Curtis, R. T.; Norton, S. P.; Parker, R. A.; Wilson, R. A. \emph{Atlas of finite groups. Maximal subgroups and ordinary characters for simple groups.} With computational assistance from J. G. Thackray. Oxford University Press, Eynsham, 1985. xxxiv+252 pp.
\bibitem[11]{Da} Dabbaghian-Abdoly, Vahid \emph{Characters of some finite groups of Lie type with a restriction containing a linear character once}. J. Algebra 309 (2007), no. 2, 543--558.
\bibitem[12]{Dor} Dornhoff, L (1971) \emph{Group Representation Theory. Part A: Ordinary Representation Theory}. Pure and Applied Mathematics. Vol. 7. New York: Marcel Dekker, Inc., pp. vii+254. [Google Scholar]
\bibitem[13]{speven} Enomoto, Hikoe \emph{The characters of the finite symplectic group} ${\rm Sp}(4,\,q)$, $q=2\sp{f}$. Osaka Math. J. \textbf{9} (1972), 75--94. MR0302750
\bibitem[14]{GH} Gardiner, Jordan C.; Humphries, Stephen P. \emph{Strong Gelfand pairs of ${\rm SL}(2,p^n)$}. Comm. Algebra 52 (2024), no. 8, 3269--3281.
\bibitem[15]{Tot2} Humphries, Stephen; Kennedy, Chelsea; Rode, Emma \emph{The total character of a finite group}. Algebra Colloq. 22 (2015), Special Issue no. 1, 775-778. (Reviewer: Silvio Dolfi)
\bibitem[16]{JaL} James, Gordon, and Liebeck, Martin. \emph{Representations and Characters of Groups}. 2nd ed., Cambridge University Press, 2001.
\bibitem[17]{JKar} Karlof, John \emph{The Subclass Algebra Associated with a Finite Group and Subgroup}, Amer. Math. Soc. Vol 207 (Jun., 1975) pp. 329-341.
\bibitem[18]{Tot1} Prajapati, S. K.; Sarma, R. \emph{Total character of a group $G$ with $(G, Z(G))$ as a generalized Camina pair}. Canad. Math. Bull. 59 (2016), no. 2, 392--402.
\bibitem[19]{Tot3} Prajapati, S. K.; Sury, B. \emph{On the total character of finite groups.} Int. J. Group Theory 3 (2014), no. 3, 47-67.
\bibitem[21]{Den} Travis, Dennis \emph{Spherical Functions on Finite Groups}, Journal of Algebra 29, 65-76 (1974).
\bibitem[22]{Wik} Wikipedia contributors. ``Gelfand pair." Wikipedia, The Free Encyclopedia, 25 May 2021.
\end{thebibliography}
\end{document}
